% !TEX root = ../main_ext_2.tex

% INTRODUCTION AARON

\section{Introduction}
Impulse responses to exogenous shocks are estimated through two main approaches, 
vector autoregressive (VAR) and Local Projection (LP) models. Since the introduction of the 
latter, an extensive literature has compared the two. One strand establishes that, under 
relatively general conditions, LPs and VARs estimate the same population impulse responses 
\parencite{plagborg2021local, fusari2026more}. In finite samples at equal lag lengths, 
however, LPs were often viewed as less sensitive to misspecification 
\parencite{jorda2005estimation, ramey2016macroeconomic}, a property that 
\textcite{olea2024double} subsequently formalise, whereas VARs were seen as more 
statistically efficient \parencite{li2024local, olea2025local}. \textcite{ludwig2026local} 
explains this finite-sample bias-variance tradeoff analytically, showing that an LP can be 
represented as a sequence of VAR estimations of increasing lag length as the horizon grows, 
and therefore argues for comparing LPs and VARs by effective complexity rather than lag length.

Our paper asks two questions. First, under correct specification, where does a fixed LP 
estimator sit relative to the full complexity-adjusted VAR frontier, and does the conventional 
equal-lag ranking survive once LP(4) is compared against the entire 
VAR-order frontier rather than only the equal-lag VAR(4)? To answer this, we benchmark a 
fixed LP(4) against progressively higher-order VAR specifications \parencite{de2025long} across 
empirically relevant sample sizes \parencite[see e.g.][]{herbst2024bias}. Second, once 
controlled functional misspecification is introduced through a nonlinear DGP, does any 
genuine LP advantage emerge, and does it manifest in point estimation, in inference, or 
both? Since \textcite{plagborg2021local} establish that LP and VAR target the same best 
linear approximation in population regardless of whether the DGP is linear, we augment the 
baseline VAR(4) with a mean-zero quadratic term in lagged state variables and examine how 
LP(4), the equal-lag VAR(4), and higher-order VARs respond as the nonlinearity strengthens.

Our findings suggest that under correct specification the conventional point-estimation 
advantage of equal-lag VARs over the LP persists, but it is largely an artifact of mismatched 
model complexity. A fixed $LP(4)$ already operates like a highly parameterized system, 
sitting high on the complexity-adjusted VAR frontier (beyond $q \approx 15$), while that 
parametric efficiency leaves VAR inference fragile and sharply under-covering near the 
unit-root boundary, where local projections retain near-nominal coverage. Introducing 
functional misspecification through a nonlinear lag-quadratic DGP leaves the shared best 
linear estimand almost unchanged, so the equal-lag VAR(4) keeps its mean squared error 
edge over LP(4) and the finite-sample gap is driven purely by variance. The nonlinearity 
instead surfaces as conditional heteroskedasticity in the residuals, which homoskedastic 
VAR delta-method intervals cannot accommodate but LP's heteroskedasticity-robust intervals 
can. Local projections therefore deliver the only valid inference, the first setting in 
this study where their robustness yields a tangible advantage.

The remainder of the paper is structured as follows. Section 2 reviews vector autoregressions 
and local projections, the empirical literature on their bias-variance trade-off at equal lag 
lengths, and their asymptotic and finite-sample equivalence. Section 3 presents the baseline 
simulation study under correct specification, Section 4 introduces functional misspecification 
through a nonlinear DGP, and Section 5 concludes.
\FloatBarrier
